\newpage
\phantomsection
\label{prf:equal-derivatives-constant-difference}

\begin{remark*}[Return]
\hyperref[cor:equal-derivatives-constant-difference]{Return to Corollary}
\end{remark*}

\begin{theorem*}[Equal Derivatives Imply Difference Is Constant]
Let $I := [a,b] \subseteq \mathbb{R}$, and let $f, g : I \to \mathbb{R}$ be continuous on $I$ and differentiable on $(a,b)$. If $f'(x) = g'(x)$ for all $x \in (a,b)$, then there exists $C \in \mathbb{R}$ such that $f(x) = g(x) + C$ for all $x \in I$.
\end{theorem*}

\begin{proof}
\textbf{Professional Standard Proof.}~\\
Define $h := f - g$. Then $h$ is continuous on $I$ and differentiable on $(a,b)$, with $h'(x) = f'(x) - g'(x) = 0$ for all $x \in (a,b)$. By the Zero Derivative Implies Constant theorem, $h$ is constant on $I$: there exists $C \in \mathbb{R}$ with $h(x) = C$ for all $x \in I$. Therefore $f(x) = g(x) + C$.
\end{proof}

\begin{proof}
\textbf{Detailed Learning Proof.}~\\

\textbf{Step 1.} Define the auxiliary function. Set $h := f - g$. Since $f$ and $g$ are both continuous on $I$ and differentiable on $(a,b)$, so is $h$.

\textbf{Step 2.} Compute the derivative of $h$. For every $x \in (a,b)$, the Sum Rule gives $h'(x) = f'(x) - g'(x) = 0$ by hypothesis.

\textbf{Step 3.} Apply the Zero Derivative Implies Constant theorem. Since $h$ is continuous on $I$, differentiable on $(a,b)$, and $h' \equiv 0$ on $(a,b)$, there exists $C \in \mathbb{R}$ such that $h(x) = C$ for all $x \in I$.

\textbf{Step 4.} Recover the conclusion. From $h(x) = C$, we have $f(x) - g(x) = C$, so $f(x) = g(x) + C$ for all $x \in I$.
\end{proof}

\begin{remark*}[Proof structure]
A one-line reduction to the preceding theorem via the auxiliary function $h = f - g$. The hypothesis $f' = g'$ becomes $h' = 0$, and the conclusion $f = g + C$ becomes $h = C$.
\end{remark*}

\begin{remark*}[Dependencies]~\\
\begin{itemize}
  \item \hyperref[thm:zero-derivative-implies-constant]{Theorem: Zero Derivative Implies Constant}
\end{itemize}
\end{remark*}
